\documentclass{amsart}
% For dealing with references we use the comment environment
\usepackage{verbatim}
\newenvironment{reference}{\comment}{\endcomment}
%\newenvironment{reference}{}{}
\newenvironment{slogan}{\comment}{\endcomment}
\newenvironment{history}{\comment}{\endcomment}

% For commutative diagrams we use Xy-pic
\usepackage[all]{xy}

% We use 2cell for 2-commutative diagrams.
\xyoption{2cell}
\UseAllTwocells

% We use multicol for the list of chapters between chapters
\usepackage{multicol}

% This is generally recommended for better output
\usepackage{lmodern}
\usepackage[T1]{fontenc}

% For cross-file-references

% Package for hypertext links:
\usepackage{hyperref}

% For any local file, say "hello.tex" you want to link to please

% Theorem environments.
%
\theoremstyle{plain}
\newtheorem{theorem}[subsection]{Theorem}
\newtheorem{proposition}[subsection]{Proposition}
\newtheorem{lemma}[subsection]{Lemma}

\theoremstyle{definition}
\newtheorem{definition}[subsection]{Definition}
\newtheorem{example}[subsection]{Example}
\newtheorem{exercise}[subsection]{Exercise}
\newtheorem{situation}[subsection]{Situation}

\theoremstyle{remark}
\newtheorem{remark}[subsection]{Remark}
\newtheorem{remarks}[subsection]{Remarks}

\numberwithin{equation}{subsection}

% Macros
%
\def\lim{\mathop{\mathrm{lim}}\nolimits}
\def\colim{\mathop{\mathrm{colim}}\nolimits}
\def\Spec{\mathop{\mathrm{Spec}}}
\def\Hom{\mathop{\mathrm{Hom}}\nolimits}
\def\Ext{\mathop{\mathrm{Ext}}\nolimits}
\def\SheafHom{\mathop{\mathcal{H}\!\mathit{om}}\nolimits}
\def\SheafExt{\mathop{\mathcal{E}\!\mathit{xt}}\nolimits}
\def\Sch{\mathit{Sch}}
\def\Mor{\mathop{\mathrm{Mor}}\nolimits}
\def\Ob{\mathop{\mathrm{Ob}}\nolimits}
\def\Sh{\mathop{\mathit{Sh}}\nolimits}
\def\NL{\mathop{N\!L}\nolimits}
\def\CH{\mathop{\mathrm{CH}}\nolimits}
\def\proetale{{pro\text{-}\acute{e}tale}}
\def\etale{{\acute{e}tale}}
\def\QCoh{\mathit{QCoh}}
\def\Ker{\mathop{\mathrm{Ker}}}
\def\Im{\mathop{\mathrm{Im}}}
\def\Coker{\mathop{\mathrm{Coker}}}
\def\Coim{\mathop{\mathrm{Coim}}}

% Boxtimes
%
\DeclareMathSymbol{\boxtimes}{\mathbin}{AMSa}{"02}

%
% Macros for moduli stacks/spaces
%
\def\QCohstack{\mathcal{QC}\!\mathit{oh}}
\def\Cohstack{\mathcal{C}\!\mathit{oh}}
\def\Spacesstack{\mathcal{S}\!\mathit{paces}}
\def\Quotfunctor{\mathrm{Quot}}
\def\Hilbfunctor{\mathrm{Hilb}}
\def\Curvesstack{\mathcal{C}\!\mathit{urves}}
\def\Polarizedstack{\mathcal{P}\!\mathit{olarized}}
\def\Complexesstack{\mathcal{C}\!\mathit{omplexes}}
% \Pic is the operator that assigns to X its picard group, usage \Pic(X)
% \Picardstack_{X/B} denotes the Picard stack of X over B
% \Picardfunctor_{X/B} denotes the Picard functor of X over B
\def\Pic{\mathop{\mathrm{Pic}}\nolimits}
\def\Picardstack{\mathcal{P}\!\mathit{ic}}
\def\Picardfunctor{\mathrm{Pic}}
\def\Deformationcategory{\mathcal{D}\!\mathit{ef}}

\setlength{\emergencystretch}{3em}
\begin{document}
\title[Stacks Project, Tag 0FUE]{466 --- Closed immersions between smooth schemes are etale-locally standard coordinate embeddings}
\maketitle
\noindent Copyright (C) 2005--2025 Johan de Jong. Stacks Project authors. Source: \href{https://stacks.math.columbia.edu/tag/0FUE}{Tag 0FUE}.

\noindent Editorial difficulty note (not author text). Difficulty H1: The construction must simultaneously lift etale coordinates on the closed subscheme and keep the defining regular sequence as transverse coordinates. The resulting map is shown etale by a fibrewise miracle-flatness argument that matches local dimensions and controls the finite residual fibre..

\noindent Editorial provenance note (not author text). History: excerpt from revision \texttt{a0444\allowbreak 6e57e\allowbreak c1fbc\allowbreak 252a8\allowbreak 71afc\allowbreak ec775\allowbreak 2fb28\allowbreak 07b14}, more-morphisms.tex, lines 4860--4937. Reference presentation was adapted; original-author context and core support blocks were retained or added. The mathematical wording is unchanged. Licensed under GNU FDL 1.2 or later; no invariant sections or cover texts. See ../licenses/COPYING. Context blocks reproduce author definitions and supporting results; they are not additional counted problems.

\begin{lemma}
\label{lemma-etale-local-structure}
Let $S$ be a scheme. Let $Y \to X$ be a closed immersion of schemes
smooth over $S$. For every $y \in Y$ there exist integers
$0 \leq m, n$ and a commutative diagram
$$
\xymatrix{
Y \ar[d] &
V \ar[l] \ar[d] \ar[r] &
\mathbf{A}^m_S
\ar[d]^{(a_1, \ldots, a_m) \mapsto (a_1, \ldots, a_m, 0 \ldots, 0)} \\
X &
U \ar[l] \ar[r]^-\pi &
\mathbf{A}^{m + n}_S
}
$$
where $U \subset X$ is open, $V = Y \cap U$,
$\pi$ is \'etale, $V = \pi^{-1}(\mathbf{A}^m_S)$, and $y \in V$.
\end{lemma}

\begin{proof}
The question is local on $X$ hence we may replace $X$ by
an open neighbourhood of $y$. Since $Y \to X$ is a regular immersion
by Divisors, Lemma
\href{https://stacks.math.columbia.edu/tag/067U}{lemma immersion smooth into smooth regular immersion (Tag 067U)}
we may assume $X = \Spec(A)$ is affine and there exists a regular sequence
$f_1, \ldots, f_n \in A$ such that $Y = V(f_1, \ldots, f_n)$.
After shrinking $X$ (and hence $Y$) further
we may assume there exists an \'etale morphism $Y \to \mathbf{A}^m_S$, see
Morphisms, Lemma \href{https://stacks.math.columbia.edu/tag/054L}{lemma smooth etale over affine space (Tag 054L)}.
Let $\overline{g}_1, \ldots, \overline{g}_m$ in $\mathcal{O}_Y(Y)$
be the coordinate functions of this \'etale morphism.
Choose lifts $g_1, \ldots, g_m \in A$ of these functions
and consider the morphism
$$
(g_1, \ldots, g_m, f_1, \ldots, f_n) :
X
\longrightarrow
\mathbf{A}^{m + n}_S
$$
over $S$. This is a morphism of schemes locally of finite presentation
over $S$ and hence is locally of finite presentation
(Morphisms, Lemma \href{https://stacks.math.columbia.edu/tag/02FV}{lemma finite presentation permanence (Tag 02FV)}).
The restriction of this morphism to
$\mathbf{A}^m_S \subset \mathbf{A}^{m + n}_S$
is \'etale by construction. Thus, in order to show that
$X \to \mathbf{A}^{m + n}_S$ is \'etale at $y$
it suffices to show that $X \to \mathbf{A}^{m + n}_S$ is flat at $y$,
see Morphisms, Lemma \href{https://stacks.math.columbia.edu/tag/02GU}{lemma etale at point (Tag 02GU)}.
Let $s \in S$ be the image of $y$. It suffices to check that
$X_s \to \mathbf{A}^{m + n}_s$ is flat at $y$, see
Theorem \href{https://stacks.math.columbia.edu/tag/039C}{theorem criterion flatness fibre (Tag 039C)}.
Let $z \in \mathbf{A}^{m + n}_s$ be the image of $y$.
The local ring map
$$
\mathcal{O}_{\mathbf{A}^{m + n}_s, z}
\longrightarrow
\mathcal{O}_{X_s, y}
$$
is flat by Algebra, Lemma \href{https://stacks.math.columbia.edu/tag/00R4}{lemma CM over regular flat (Tag 00R4)}.
Namely, schemes smooth over fields are regular and regular rings
are Cohen-Macaulay, see Varieties, Lemma \href{https://stacks.math.columbia.edu/tag/056S}{lemma smooth regular (Tag 056S)}
and Algebra, Lemma \href{https://stacks.math.columbia.edu/tag/00NQ}{lemma regular ring CM (Tag 00NQ)}.
Thus both source and target are regular local rings (and hence CM).
The source and target have the same dimension: namely, we have
$\dim(\mathcal{O}_{Y_s, y}) = \dim(\mathcal{O}_{\mathbf{A}^m_s, z})$
by More on Algebra, Lemma \href{https://stacks.math.columbia.edu/tag/07QP}{lemma dimension etale extension (Tag 07QP)},
we have $\dim(\mathcal{O}_{\mathbf{A}^{m + n}_s, z}) = n +
\dim(\mathcal{O}_{\mathbf{A}^m_s, z})$, and we have
$\dim(\mathcal{O}_{X_s, y}) = n + \dim(\mathcal{O}_{Y_s, y})$
because $\mathcal{O}_{Y_s, y}$ is the quotient of
$\mathcal{O}_{X_s, y}$ by the regular sequence $f_1, \ldots, f_n$
of length $n$ (see
Divisors, Remark \href{https://stacks.math.columbia.edu/tag/0FUD}{remark relative regular immersion elements (Tag 0FUD)}).
Finally, the fibre ring of the displayed arrow is finite over $\kappa(z)$
since $Y_s \to \mathbf{A}^m_s$ is \'etale at $y$.
This finishes the proof.
\end{proof}
\end{document}
